%------------------------------------------------------------------------------%
\begin{question}
Use escalonamento para analisar o sistema
\[
  \systeme{
   x +  y = 4,
  2x + 2y = 8
  }
\]
\end{question}

%------------------------------------------------------------------------------%
\begin{resolution}{Escalonamento}
\begin{columns}
\column{0.4\linewidth}
  \[\;\;\,
    \left[
    \begin{array}{cc|c}
      1 & 1 & 4 \\
      2 & 2 & 8
    \end{array}
    \right]
  \]
  \pause
\column{0.6\linewidth}
  Aplicamos
  \\ \pause
  \(
    L_2 \leftarrow L_2 - 2 L_1
  \)
\end{columns}
\vspace{-\baselineskip}
\[
  \pause
  L'_2
  \;=\;
  L_2 - 2 L_1
  \pause
  \;=\;
  \left[
  \begin{array}{cc|c}
    2 & 2 & 8
  \end{array}
  \right]
  - 2
  \left[
  \begin{array}{cc|c}
    1 & 1 & 4
  \end{array}
  \right]
  \pause
  \;=\;
  \left[
  \begin{array}{cc|c}
    0 & 0 & 0
  \end{array}
  \right]
\]

\pause
\[
\left[
\begin{array}{cc|c}
  1 & 1 & 4 \\
  0 & 0 & 0
\end{array}
\right]
\]
\end{resolution}

%------------------------------------------------------------------------------%
\begin{resolution}{Solução}
A segunda linha representa
\[
0 = 0
\]

\pause
\bigskip
Portanto, o sistema possui \emph{infinitas} soluções
\end{resolution}

%------------------------------------------------------------------------------%
\endinput
